3.6. 设 $V$ 为 $S$ 上秩为 $n+2$ 的局部自由向量丛, $X$ 为 $V$ 上普通二次型 $Q$ 定义的 $S$ 上光滑二次超曲面. 再设 $H$ 为 $\mathbb P(V^*)$ 中与 $X$ 横截相交的超平面, 置 $X^o=X-H$. 概形 $X$ 和 $Y=X\cap H$ 是 $S$ 上维数分别为 $n$ 和 $n-1$ 的光滑二次超曲面. 因而可以用 3.3 计算 $X^o=X-Y$ 的上同调.

%%%%%%%%%% source scan idx 86 | folio 79 %%%%%%%%%%

\begin{equation}\tag{3.6.1}
\begin{tikzcd}[column sep=huge,row sep=huge]
X^o \arrow[r,hook] \arrow[rd,"f"']
  & X \arrow[d,"p"]
  & Y \arrow[l,hook] \arrow[ld,"q"]\\
{} & S & \mathord{.}
\end{tikzcd}
\end{equation}

若 $n=0$, 则 $Y=\varnothing$ 且 $X=X^o$. 假设 $n>0$. 有两个互为对偶的长正合列:
\begin{equation}\tag{3.6.2}
\cdots\xrightarrow{\partial}R^if_!\mathbb Z_\ell
\longrightarrow R^ip_*\mathbb Z_\ell
\xrightarrow{r_i}R^iq_*\mathbb Z_\ell\longrightarrow\cdots
\end{equation}
以及
\begin{equation}\tag{3.6.3}
\cdots\xrightarrow{\partial}R^{i-2}q_*\mathbb Z_\ell(-1)
\longrightarrow R^ip_*\mathbb Z_\ell
\longrightarrow R^if_*\mathbb Z_\ell\longrightarrow\cdots.
\end{equation}

当 $n$ 为偶数时, 限制同态 $r_i$ 在 $i\neq n,2n$ 时为同构; 当 $n$ 为奇数时, 它在 $i\neq n-1,2n$ 时为同构.\EditorialFootnote{法文原印本与英文底本在奇数维情形均漏列例外 $i=2n$. 此时 $X$ 的顶次上同调非零, 而 $Y$ 的该次上同调为零; 下文关于紧支集上同调的结论也要求此例外.} 若 $n=2m$ 且 $\alpha$ 为一个母空间的上同调类, 则
\begin{equation*}
r_n(\alpha)=\tfrac12\eta^m,
\end{equation*}
所以 $r_n$ 为满射, 其核为 $R^np_*\mathbb Z_\ell$ 的本原部分. 若 $n=2m+1$, 则 $r_{2m}(\eta^m)=\eta^m$, 因而 $r_{2m}$ 为单射, 其余核为 $R^{2m}q_*\mathbb Z_\ell$ 的本原商. 由此得到当 $i\neq n,2n$ 时 $R^if_!\mathbb Z_\ell=0$, 而当 $i=n,2n$ 时 $R^if_!\mathbb Z_\ell$ 为秩一的扭常自由 $\mathbb Z_\ell$-层. 由对偶性, 或由 (3.6.3), 同样可证当 $i\neq0,n$ 时 $R^if_*\mathbb Z_\ell=0$, 且 $f_*\mathbb Z_\ell=\mathbb Z_\ell$, 而 $R^nf_*\mathbb Z_\ell$ 的秩为一.

\footnote{印本这一段将 $n=2m+1$ 写作 $r=2m+1$, 并将所列正合列 (3.6.3) 引用为 3.5.3.}

假设 $n=2m$, 保留 3.3 的记号.

%%%%%%%%%% source scan idx 87 | folio 80 %%%%%%%%%%

\begin{center}
\begin{tikzcd}[column sep=large,row sep=large]
0 \arrow[r]
 & R^nf_!\mathbb Z_\ell(m) \arrow[r] \arrow[d,"\varphi"]
 & R^np_*\mathbb Z_\ell(m) \arrow[r] \arrow[d,equal]
 & R^nq_*\mathbb Z_\ell(m) \arrow[r]
 & 0\\
0
 & R^nf_*\mathbb Z_\ell(m) \arrow[l]
 & R^np_*\mathbb Z_\ell(m) \arrow[l]
 & R^{n-2}q_*\mathbb Z_\ell(m-1) \arrow[l]
 & 0\mathord{.} \arrow[l]
\end{tikzcd}
\end{center}

此时层 $R^nf_!\mathbb Z_\ell(m)$ 有两个自然的互为相反数的生成元 $\pm\delta$, 它们在 $R^np_*\mathbb Z_\ell(m)$ 中的像为 $\pm(c\ell(\alpha)-c\ell(\beta))$. 由 3.3(iii)(b),
\begin{equation*}
\operatorname{Tr}(\delta^2)
=(c\ell(\alpha)-c\ell(\beta))^2=(-1)^m\cdot2,
\end{equation*}
且 $\pm\varphi(\delta)$ 是 $R^nf_*\mathbb Z_\ell(m)$ 的自然对偶生成元 $\pm\delta'$ 的两倍, 其中这两个对偶生成元即为 $c\ell(\alpha)$ 和 $c\ell(\beta)$ 的互为相反数的像.

现在假设 $n=2m+1$, 对 $Y$ 使用 3.3 的记号.

\begin{center}
\begin{tikzcd}[column sep=large,row sep=large]
0 \arrow[r]
 & R^{2m}p_*\mathbb Z_\ell(m) \arrow[r]
 & R^{2m}q_*\mathbb Z_\ell(m) \arrow[r,"\partial"]
 & R^nf_!\mathbb Z_\ell(m) \arrow[r] \arrow[d,"\varphi"]
 & 0\\
{}
 & R^{2m+2}p_*\mathbb Z_\ell(m+1)
 & R^{2m}q_*\mathbb Z_\ell(m) \arrow[l]
 & R^nf_*\mathbb Z_\ell(m+1) \arrow[l]
 & 0\mathord{.} \arrow[l]
\end{tikzcd}
\end{center}

\footnote{印本下行最左项的扭为 $(m)$; 对偶性及其余三项要求这里给出的扭 $(m+1)$.}

此时层 $R^nf_!\mathbb Z_\ell(m)$ 有两个自然的互为相反数的生成元 $\pm\delta$, 即 $\partial c\ell(\alpha)$ 和 $\partial c\ell(\beta)$. 有
\begin{equation*}
\delta\,\varphi(\delta)=\delta^2
=(\partial c\ell(\alpha))^2
=\partial\bigl(c\ell(\alpha)\,\partial c\ell(\alpha)\bigr)=0,
\end{equation*}
因而 $\varphi(\delta)=0$. 另一方面, $R^nf_*\mathbb Z_\ell(m+1)$ 有两个自然的互为相反数的生成元 $\pm\delta'$, 它们在 $R^{2m}q_*\mathbb Z_\ell(m)$ 中的像为 $\pm(c\ell(\alpha)-c\ell(\beta))$, 且 $\operatorname{Tr}(\delta\delta')=\pm1$.

%%%%%%%%%% source scan idx 88 | folio 81 %%%%%%%%%%

上述结果汇集于下表.

\underline{表} 3.7. \underline{维数为 $n$ 的仿射二次超曲面 $X^o=X-Y$ 的上同调.}

\medskip
\noindent\textbf{A.} 上同调无挠. 非零的 Betti 数为
\begin{align*}
\text{普通上同调:}\quad &b_0=b_n=1,\\
\text{紧支集上同调:}\quad &b_{2n}=b_n=1,
\end{align*}
唯当 $n=0$ 时, 应为 $b_0=2$.

\medskip
\noindent\textbf{B.} 若 $n=2m>0$, 则
\begin{align*}
R^nf_!\mathbb Z_\ell(m)
  &=\text{层 }R^{2m}p_*\mathbb Z_\ell(m)\text{ 的本原部分},\\
R^nf_*\mathbb Z_\ell(m)
  &=\text{层 }R^{2m}p_*\mathbb Z_\ell(m)\text{ 的本原商}.
\end{align*}
若 $n=2m+1$, 则
\begin{align*}
R^nf_!\mathbb Z_\ell(m)
  &=\text{层 }R^{2m}q_*\mathbb Z_\ell(m)\text{ 的本原商},\\
R^nf_*\mathbb Z_\ell(m+1)
  &=\text{层 }R^{2m}q_*\mathbb Z_\ell(m)\text{ 的本原部分}.
\end{align*}
局部地, 这些群有精确到符号的自然生成元; 以 $\delta$ 表示紧支集上同调的这种生成元, 以 $\delta'$ 表示普通上同调的这种生成元.

\medskip
\noindent\textbf{C.} 映射
\begin{equation*}
\varphi:R^nf_!\mathbb Z_\ell\longrightarrow R^nf_*\mathbb Z_\ell
\end{equation*}
在 $n$ 为奇数时为零, 在 $n>0$ 为偶数时将 $\pm\delta$ 映为 $\pm2\delta'$. 此外,
\begin{equation*}
\operatorname{Tr}(\delta\delta')=\pm1,
\qquad
\operatorname{Tr}(\delta^2)=
\begin{cases}
0, & n=2m+1,\\
(-1)^m\cdot2, & n=2m.
\end{cases}
\end{equation*}

\underline{验证} 3.8. $\mathbb C^n$ 中方程为 $\sum_i z_i^2=1$ 的复仿射二次超曲面 $S$ 微分同胚于实球面 $S\cap\mathbb R^n$ 的切丛. 这与 A 和 C 相符.

\medskip
\underline{参考文献}

\begin{thebibliography}{9}
\bibitem{MumfordGIT}
D. Mumford, \emph{Geometric Invariant Theory}, Ergebnisse \textbf{34},
Springer-Verlag, 1965.
\end{thebibliography}
